% Discussão da campanha oficial integralmente auditada.
\subsection{Compromissos observados entre custo e qualidade}

Na amostra obtida, o $\varepsilon$-restrito alcançou o menor custo,
R\$~\num{60460000}, enquanto a soma ponderada atingiu os menores desvios
de $SiO_2$ e $Al_2O_3$, respectivamente \num{2.406656} e \num{1.452043}
p.p.$^2$. Esses extremos não constituem uma única solução ideal. As
alternativas completas que os originaram estão na
Tabela~\ref{tab:tp2_extremos}, com composições e proveniência recuperáveis
no conjunto de dados.

Por exemplo, a alternativa de menor custo do $\varepsilon$-restrito
apresenta o vetor, com custo em milhões de reais,
$(\num{60.460},\num{6.343},\num{5.706})$. A alternativa ponderada de menor
desvio de $SiO_2$ tem vetor
$(\num{63.700},\num{2.407},\num{3.284})$.
Nesse par, um acréscimo de R\$~\num{3240000} reduz ambos os desvios,
evidenciando o compromisso entre economia e qualidade.

A alternativa ponderada de menor desvio de $Al_2O_3$ tem vetor
$(\num{65.660},\num{5.498},\num{1.452})$.
Em relação à alternativa ponderada de menor desvio de $SiO_2$, custa
R\$~\num{1960000} a mais, melhora $Al_2O_3$ e piora $SiO_2$.
Portanto, o conflito também ocorre entre as duas medidas de qualidade,
não apenas entre custo e um indicador agregado.

O mínimo individual corrigido do TP1 para custo foi R\$~\num{60300000}.
Para os desvios de $SiO_2$ e $Al_2O_3$, os mínimos foram, respectivamente,
\num{2.372757} e \num{1.436169} p.p.$^2$. A coleta multiobjetivo não
superou esses três mínimos marginais; sua finalidade é explorar
compromissos, e não repetir a otimização isolada de cada objetivo.
Também foi encontrado $f_2$ de \num{6.872459} p.p.$^2$ na união ponderada,
acima do máximo empírico TP1 de \num{6.705745}; o valor normalizado é
aproximadamente \num{1.038}. Isso é permitido pela normalização sem
truncamento e não implica violar as restrições físicas.

\subsection{Semelhança e variabilidade entre as cinco varreduras}

As cinco varreduras ponderadas encontraram intervalos marginais bastante
semelhantes. Entre os dez pares de varreduras, a menor razão entre o
comprimento da interseção e o da união dos intervalos foi
\num{91.53}\% para custo, \num{91.42}\% para $SiO_2$ e \num{87.47}\%
para $Al_2O_3$. No $\varepsilon$-restrito, esses mínimos foram
\num{67.49}\%, \num{87.32}\% e \num{58.39}\%, respectivamente.
A medida é apenas descritiva e marginal: não demonstra que as composições
ou a distribuição dos pontos em 3D sejam iguais.

Nas seis projeções das Figuras~\ref{fig:tp2_ponderada_1_2}
a~\ref{fig:tp2_epsilon_2_3}, as cores e símbolos das cinco repetições
ocupam faixas amplamente sobrepostas dentro de cada abordagem. A projeção
custo versus $Al_2O_3$ mostra bandas decrescentes comuns às repetições;
na projeção das duas qualidades, observam-se agrupamentos e lacunas,
sobretudo no $\varepsilon$-restrito. Assim, as repetições alcançaram regiões
projetadas semelhantes, com diferenças de extensão e densidade local,
sem evidência de identidade entre os conjuntos tridimensionais.

Os arquivos por repetição contêm de \num{1617} a \num{1963} composições
ponderadas e de \num{634} a \num{790} composições $\varepsilon$-restritas
ND. Assim, contagens diferentes coexistem com intervalos sobrepostos;
uma varredura com mais composições não necessariamente amplia a região
de objetivos alcançada. As projeções identificam a distribuição de cada
repetição; a união por método preserva a proveniência e filtra novamente
em três objetivos.

Em todas as sementes, as 16 configurações ponderadas encontraram viável,
enquanto nove dos 16 pares $\varepsilon$ encontraram viável. As 35 ausências
se concentram nos sete pares com pelo menos um limite normalizado zero:
$(0,0)$, $(0,1/3)$, $(0,2/3)$, $(0,1)$, $(1/3,0)$, $(2/3,0)$ e $(1,0)$.
Esse padrão sugere dificuldade persistente da busca em alcançar os
extremos empíricos de qualidade sob os limites simultâneos. Não demonstra
que esses subproblemas sejam matematicamente inviáveis, nem autoriza
substituir as sementes por repetições bem-sucedidas.

\subsection{Limitações e implicações para a continuidade}

A soma ponderada produziu maior cobertura descritiva da amostra oposta,
conforme a comparação bidirecional em três objetivos, e alcançou extremos
de qualidade mais próximos dos resultados mono-objetivo. O
$\varepsilon$-restrito forneceu o menor custo e alternativas que não são
todas dominadas pela união ponderada. As duas amostras contribuem para
explicitar preferências diferentes; cardinalidade, mínimos marginais e
cobertura devem ser considerados em conjunto.

As conclusões se restringem à instância, ao GVNS, às duas malhas e ao teto
de \num{10000} avaliações por configuração. Malhas finitas podem omitir
regiões relevantes; a escalarização ponderada não recupera necessariamente
toda a fronteira de um problema discreto. O orçamento e as sementes
limitam a exploração, e o arquivo contém apenas soluções observadas pela
heurística. Não se dispõe de prova de otimalidade ou de completude da
fronteira. A comparação de tempos é descritiva, com custo adicional de
observação, persistência, consolidação e auditoria explicitado no protocolo.

\section{Considerações Finais da Entrega 2}

Foram integradas duas formulações escalares normalizadas ao GVNS existente,
com orçamento comum, cinco varreduras por abordagem, preservação da melhor
viável e dos candidatos não dominados, registros completos e verificação
independente. As seis projeções e as tabelas permitem examinar custo,
qualidades e variabilidade sem substituir a comparação tridimensional por
dominância aparente em duas dimensões.

A escolha de uma ação final permanece para a etapa multicritério seguinte.
As composições e suas origens foram preservadas para permitir recuperar os
três objetivos e calcular atributos adicionais sem reconstruir receitas
a partir de gráficos ou valores arredondados.
